%! Function Model Selection and Construction — Unit 1, Lesson 1.7 — Slides (Beamer)
\documentclass[aspectratio=169, 11pt]{beamer}
\usepackage{precalculus-beamer}   % provides \plumheader{} and \sectionlabel[color]{}

\begin{document}

% TITLE SLIDE
{
\setbeamercolor{background canvas}{bg=plum}
\begin{frame}[plain]
  \vspace{2.8cm}\hspace{0.55cm}%
  \begin{minipage}{0.9\textwidth}
    {\color{goldacc}\bfseries\small LESSON 1.7}\\[0.5cm]
    {\color{white}\bfseries\LARGE Function Model Selection and Construction}\\[1.2cm]
    {\color{lilacmid}\small Precalculus~$\cdot$~Unit 1: Functions and Change}
  \end{minipage}
\end{frame}
}

% HOOK SLIDE
\begin{frame}[t, plain]
  \plumheader{Hook: Eighty Feet of Fence and a Barn}
  \sectionlabel[goldacc]{Vote First, Algebra Second}

  The garden club has exactly \textbf{80 feet of fencing} for a rectangular pen.
  One whole side runs along the barn wall --- that side needs \textbf{no fence at
  all}.

  \bigskip
  \begin{columns}[T]
    \column{0.46\textwidth}
      \begin{block}{Question 1}
        Should the pen be a \textbf{square}? Vote now, before anyone computes
        anything.
      \end{block}
    \column{0.46\textwidth}
      \begin{block}{Question 2}
        What \textbf{kind} of function turns a width into an area --- and how
        would you know before you built it?
      \end{block}
  \end{columns}

  \bigskip
  \textit{Today we answer that twice: from a table of measurements, and from the
  fencing constraint itself. The two answers had better agree.}
\end{frame}

% SECTION 1 — READING THE MODEL OFF A TABLE
\begin{frame}[t, plain]
  \plumheader{1. Reading the Model Off a Table}

  \sectionlabel[redacc]{Check the Input Row First --- Differences Need Equal Spacing}

  \begin{columns}[T]
    \column{0.52\textwidth}
      \begin{center}
      \begin{tabular}{|c|c|c|c|c|c|}
      \hline
      \(x\) & 0 & 1 & 2 & 3 & 4 \\
      \hline
      \(y\) & 5 & 8 & 11 & 14 & 17 \\
      \hline
      1st & --- & 3 & 3 & 3 & 3 \\
      \hline
      \end{tabular}

      \medskip
      \begin{tabular}{|c|c|c|c|c|c|}
      \hline
      \(w\) (ft) & 5 & 10 & 15 & 20 & 25 \\
      \hline
      \(A\) (ft\(^2\)) & 350 & 600 & 750 & 800 & 750 \\
      \hline
      1st & --- & 250 & 150 & 50 & \(-50\) \\
      \hline
      2nd & --- & --- & \(-100\) & \(-100\) & \(-100\) \\
      \hline
      \end{tabular}
      \end{center}

    \column{0.44\textwidth}
      \begin{block}{The decision tree}
        1st differences constant \(\Rightarrow\) \textbf{linear} \\[3pt]
        2nd differences constant \(\Rightarrow\) \textbf{quadratic} \\[3pt]
        3rd differences constant \(\Rightarrow\) \textbf{cubic} \\[3pt]
        \(n\)th constant \(\Rightarrow\) degree \(n\)
      \end{block}

      \medskip
      \begin{block}{What it means}
        Constant \textbf{first} differences: the \textit{amount} of change
        repeats.\\[3pt]
        Constant \textbf{second} differences: the \textit{rate} of change is
        changing steadily.
      \end{block}
  \end{columns}
\end{frame}

% SECTION 2 — READING THE MODEL OFF THE CONTEXT
\begin{frame}[t, plain]
  \plumheader{2. Reading the Model Off the Context}

  No table? The description carries the clue.

  \bigskip
  \begin{columns}[T]
    \column{0.54\textwidth}
      \begin{center}
      \begin{tabular}{|l|l|}
      \hline
      \textbf{What the context says} & \textbf{Model} \\
      \hline
      Same amount added each step & Linear \\
      \hline
      An \textbf{area} (two lengths) & Quadratic \\
      \hline
      A \textbf{volume} (three lengths) & Cubic \\
      \hline
      Rises, falls, rises again & Higher degree \\
      \hline
      \textbf{Inversely proportional} & Rational \\
      \hline
      Rule \textbf{changes at a threshold} & Piecewise \\
      \hline
      \end{tabular}
      \end{center}

    \column{0.42\textwidth}
      \begin{block}{Inverse proportion}
        \(y = \frac{k}{x^n}\) --- as \(x\) grows, \(y\) shrinks. These are
        \textbf{rational} functions; Unit 3 is about them. Today: recognize the
        words.
      \end{block}

      \medskip
      \begin{block}{A counting fact}
        A degree-\(n\) polynomial fits through \(n+1\) points with distinct
        inputs.\\[3pt]
        So ``it fits the data'' is not by itself a reason to believe a model.
      \end{block}
  \end{columns}
\end{frame}

% SECTION 3 — BUILDING THE MODEL FROM THE CONSTRAINT
\begin{frame}[t, plain]
  \plumheader{3. Building the Model from the Constraint}

  \sectionlabel[goldacc]{The Four Steps}

  \begin{columns}[T]
    \column{0.36\textwidth}
      \begin{center}
      \begin{tikzpicture}[x=0.8cm, y=0.8cm]
        \fill[lilac] (0,0) rectangle (4.4,1.7);
        \draw[very thick, slate] (-0.4,1.7) -- (4.8,1.7);
        \node[slate, above, font=\scriptsize] at (2.2,1.75) {barn wall --- no fence};
        \draw[very thick, plum] (0,1.7) -- (0,0) -- (4.4,0) -- (4.4,1.7);
        \node[plum, left,  font=\small] at (0,0.85) {\(w\)};
        \node[plum, right, font=\small] at (4.4,0.85) {\(w\)};
        \node[plum, below, font=\small] at (2.2,0) {\(\ell\)};
      \end{tikzpicture}
      \end{center}

    \column{0.30\textwidth}
      \begin{block}{Steps}
        \begin{enumerate}\setlength{\itemsep}{4pt}
          \item Name the variables.
          \item Write the \textbf{constraint}.
          \item Solve it for one variable.
          \item \textbf{Substitute} into what you want.
        \end{enumerate}
      \end{block}

    \column{0.30\textwidth}
      \begin{block}{The pen}
        \(2w + \ell = 80\) \\[3pt]
        \(\ell = 80 - 2w\) \\[3pt]
        \(A = w\ell = w(80 - 2w)\) \\[3pt]
        \(A(w) = 80w - 2w^2\), \ \(0 < w < 40\)
      \end{block}
  \end{columns}

  \bigskip
  \textit{Check: \(A(20) = 1600 - 800 = 800\) --- exactly what the table said. The
  table and the constraint produced the same function.}
\end{frame}

% SECTION 4 — USING THE MODEL
\begin{frame}[t, plain]
  \plumheader{4. Using the Model}

  \begin{columns}[T]
    \column{0.48\textwidth}
      \begin{block}{Average rate of change --- with units}
        \(\frac{A(20)-A(10)}{20-10} = \frac{800-600}{10} = 20\) \\[3pt]
        \(\frac{A(30)-A(20)}{30-20} = \frac{600-800}{10} = -20\) \\[5pt]
        Units: \textbf{square feet per foot of width}.
      \end{block}

      \medskip
      \begin{block}{The sign change}
        \(+\) then \(-\) means the largest area sits \textbf{between} the two
        intervals.
      \end{block}

    \column{0.48\textwidth}
      \begin{block}{The vertex}
        \(w = -\frac{b}{2a} = -\frac{80}{2(-2)} = 20\) \\[3pt]
        \(\ell = 80 - 2(20) = 40\) \\[3pt]
        Largest area \(= 800\) ft\(^2\).
      \end{block}

      \medskip
      \begin{block}{Back to the vote}
        The best pen is \(20 \times 40\) --- \textbf{not} a square.\\[3pt]
        One side is free, so it pays to make the free side the long one.
      \end{block}
  \end{columns}
\end{frame}

% SECTION 5 — THE FINE PRINT
\begin{frame}[t, plain]
  \plumheader{5. The Fine Print: Restrictions and Assumptions}

  \begin{columns}[T]
    \column{0.48\textwidth}
      \begin{block}{Restrictions}
        Domain: \(0 < w < 40\) --- widths outside this are not bad arithmetic,
        they are \textbf{outside the domain}.\\[3pt]
        Range: \(0 < A \le 800\) square feet.
      \end{block}

      \medskip
      \begin{block}{Assumptions}
        All 80 ft is used \(\cdot\) the pen is exactly rectangular \(\cdot\) no
        gate gap \(\cdot\) the barn wall is at least 40 ft long.
      \end{block}

    \column{0.48\textwidth}
      \begin{block}{When one rule is not enough}
        Suppose the barn is only \textbf{30 ft} long.\\[3pt]
        A pen with \(\ell > 30\) needs fence on part of the fourth side too ---
        the rule \textbf{changes} at that threshold.\\[3pt]
        Different rules on different intervals \(=\) a \textbf{piecewise} model.
      \end{block}

      \medskip
      \textit{Naming the assumption is part of the answer. A model that fits the
      numbers can still be wrong about the world.}
  \end{columns}
\end{frame}

% CLASS PRACTICE
\begin{frame}[t, plain]
  \plumheader{Class Practice}
  \sectionlabel[redacc]{Try It}

  \begin{enumerate}\setlength{\itemsep}{10pt}
    \item Take differences until they are constant, then name the model type:
          \(y = 0,\ 1,\ 8,\ 27,\ 64\) at \(x = 0,1,2,3,4\).
    \item A garden bed is built against a wall using 30 feet of edging on the
          other three sides.
          \begin{enumerate}[(a)]\setlength{\itemsep}{4pt}
            \item Write the constraint and build \(A(w)\).
            \item State the domain the context allows.
          \end{enumerate}
    \item A city charges \$0.05 per gallon for the first 5{,}000 gallons and
          \$0.09 per gallon after that. What kind of model is the water bill, and
          what is the threshold?
  \end{enumerate}
\end{frame}

\end{document}
